% ECONOMICS_PROBLEM: {"id": "C90-639", "title": "Lab–online–field population generalizability", "jel_code": "C90", "source_id": "OP-0070"}
% !TeX program = xelatex
\documentclass[11pt,a4paper]{article}
\usepackage[margin=23mm]{geometry}
\usepackage{fontspec}
\setmainfont{Latin Modern Roman}
\usepackage{amsmath,amssymb,mathtools}
\usepackage{xcolor,enumitem,booktabs,longtable,array,xurl,hyperref}
\definecolor{muted}{HTML}{53636C}
\hypersetup{colorlinks=true,urlcolor=blue,linkcolor=blue}
\setlength{\parindent}{0pt}
\setlength{\parskip}{5pt}
\newcommand{\R}{\mathbb R}
\newcommand{\E}{\mathbb E}
\newcommand{\N}{\mathbb N}
\newcommand{\ind}{\mathbf 1}
\newcommand{\Var}{\operatorname{Var}}
\newcommand{\Cov}{\operatorname{Cov}}
\newcommand{\supp}{\operatorname{supp}}
\DeclareMathOperator*{\argmin}{arg\,min}
\DeclareMathOperator*{\argmax}{arg\,max}
\newcommand{\entrymeta}[3]{{\sffamily\small\color{muted}\textbf{\texttt{#1}}\quad|\quad #2\par #3\par}}
\newcommand{\Problemref}[1]{\texttt{#1}}
\newcommand{\JELref}[2]{\texttt{#2} (JEL \texttt{#1})}
\begin{document}
\section*{Lab–online–field population generalizability}
\textbf{Problem ID:} \texttt{C90-639}\quad \textbf{JEL:} \texttt{C90}\par
% BEGIN ECONOMICS STATEMENT
\label{problem:OP-0070}
\label{atlas:prob:B10}

\noindent\textbf{Original identifier:} B10 (atlas item 70). \textbf{Classification:} Population parameter generalizability research program. \textbf{Original tractability label:} Easy (editorial, not a complexity result).

\subsection*{Definitions and assumptions}

Let populations $p\in\mathcal P$ represent laboratory, online, and field recruitment frames, all facing a common experimental protocol $G$. For person $i$, let $\theta_i\in\Theta$ be a latent behavioral parameter and $X_i$ observed pretreatment characteristics. The measurement model is $P(A_i\mid G,\theta_i,X_i,p)$, with common normalization and a specified noise distribution. Let $R_i=1$ denote participation and define the target population distribution independently of recruitment. Parameter comparability requires measurement invariance: identical latent types facing the same protocol generate identical conditional action distributions across frames, unless population-specific measurement differences are modeled.

\subsection*{Formal research question}

Identify or sharply bound $F_{\theta\mid p}$ and the extent to which between-population differences are attributable to observed composition. Under conditional latent-type exchangeability and overlap, evaluate
\[
F_{\theta\mid p_t}(z)=\int F_{\theta\mid X=x,R=1,p_s}(z)\,dF_{X\mid p_t}(x),\quad z\in\Theta,
\]
interpreted componentwise for vector parameters. Determine sensitivity to violations of sampling exchangeability and measurement invariance. Test whether parameters calibrated on one frame predict choices on held-out frames within a declared discrepancy metric and tolerance, accounting for first-stage estimation uncertainty.

\subsection*{Interpretation and scope}

This target differs from B1: it concerns transport of a latent parameter distribution, not merely transport of one treatment effect. Latent parameters may be unidentified even when an average treatment effect is experimentally identified. Subject-pool regressions should therefore account for estimation noise, rather than regress noisy individual estimates as if observed without error. Education or digital literacy can explain statistical variation without being its causal source; causal statements about these traits require a separate intervention or identification argument. Representative panels reduce some sampling discrepancies but do not eliminate participation selection or protocol differences. Common instructions, incentives, device access, and comprehension checks help support measurement invariance, although checks conducted after treatment must not create new selection bias. A benchmark infrastructure should preregister target populations and protocols, document exclusions, and evaluate both population-average and distributional predictions. Good fit in observed populations cannot establish universal generalizability to every future subject pool.

\subsection*{Source audit and status}

[S1] and [S2] verify the original generalizability citations. [S3] again concerns laboratory replicability, so its role is background quality control rather than direct evidence of population parameter transport. The source overlap with B1 is retained transparently while the formal targets are distinguished. This editorial measurement and transportability problem does not inherit a universally easy classification from the atlas; identification difficulty depends on sampling and measurement restrictions.

\subsection*{References}

\begin{enumerate}[label={[S\arabic*]},leftmargin=*]

\item Glenn W. Harrison and John A. List (2004). \emph{Field Experiments}. Journal of Economic Literature 42(4), 1009--1055. \url{https://www.aeaweb.org/articles?id=10.1257/0022051043004577}\par \textit{Source role:} Taxonomy of field settings; methodological foundation. \textit{Locator:} Abstract; six dimensions of field context.

\item Steven D. Levitt and John A. List (2007). \emph{What Do Laboratory Experiments Measuring Social Preferences Reveal About the Real World?}. Journal of Economic Perspectives 21(2), 153--174. \url{https://www.aeaweb.org/articles?id=10.1257/jep.21.2.153}\par \textit{Source role:} Methodological discussion of generalizability. \textit{Locator:} Abstract; discussion of scrutiny, context, selection, and stakes.

\item Colin F. Camerer, Anna Dreber, Eskil Forsell, Teck-Hua Ho, Jürgen Huber, Magnus Johannesson, Michael Kirchler, Johan Almenberg, Adam Altmejd, Taizan Chan, Emma Heikensten, Felix Holzmeister, Taisuke Imai, Siri Isaksson, Gideon Nave, Thomas Pfeiffer, Michael Razen, and Hang Wu (2016). \emph{Evaluating Replicability of Laboratory Experiments in Economics}. Science 351(6280), 1433--1436. \url{https://doi.org/10.1126/science.aaf0918}\par \textit{Source role:} Evidence on laboratory replicability, not field transportability. \textit{Locator:} Abstract; author repository metadata and supplemental materials.

\end{enumerate}

\subsection*{Common conventions: Source mathematical conventions}

Unless an entry states otherwise, agent and firm sets are finite. State and action spaces are measurable, strategies and outcome maps are measurable, probability laws are specified, and expectations exist for the targets under discussion. Infinite-horizon discount factors lie in $(0,1)$ and discounted objectives are finite. Norms, tolerances and complexity input sizes must be specified for computational comparisons. Constants or primitives that appear locally are inputs, not empirical facts asserted by this catalogue.

\textbf{Identification.} A statistical model $m\in\mathcal M$ implies an observable law $P_m$ and target $\theta(m)$. At law $P$, its identified set is
\[
\Theta(P;\mathcal M)=\{\theta(m):m\in\mathcal M,\ P_m=P\}.
\]
Point identification holds when this set is a singleton, or equivalently when $P_{m_1}=P_{m_2}$ implies $\theta(m_1)=\theta(m_2)$ for all compatible models. Sharp bounds mean bounds attained, or approached when the boundary is not attained, by compatible models. Writing this set is a definition, not a solution: the substantive task is to establish informative restrictions and derive or estimate the set. An unrestricted-confounding model may give only trivial bounds.

\textbf{Inference.} Identification, estimability and valid uncertainty quantification are separate questions. Fix $\alpha\in(0,1)$. A confidence procedure $C_n$ over a declared law class $\mathcal P$ has asymptotic uniform coverage at level $1-\alpha$ when
\[
\liminf_{n\to\infty}\inf_{P\in\mathcal P}\Pr_P\{\theta(P)\in C_n\}\geq1-\alpha.
\]
For set-identified targets the coverage event must be defined separately, for example coverage of the full identified set. If an entry uses identification-indexed models instead of uniquely indexed laws, the same coverage convention is applied over those models. Pointwise limits do not establish uniform validity, and asymptotic validity does not guarantee finite-sample precision. Sample sizes, dependence structures and weak-identification sequences are part of each application.

\textbf{Counterfactuals.} $Y_i(a)$ is a potential outcome under a specified intervention, including its timing, implementation and versions. It is not $Y_i$ conditional on voluntarily choosing $a$. A full intervention vector permits interference; shorthand scalar treatment notation does not silently impose no interference. Assignment, selection, attrition, anticipation and measurement must be specified. A claim about an instrument's local effect is not automatically a population policy effect.

\textbf{Economic equilibrium.} A counterfactual model includes best responses, constraints, feasibility, market clearing, state transitions and expectations consistent with the stated information. When several equilibria exist, either a declared selector is an input or the target is set-valued. Comparisons across policy regimes must use the stated selection convention. Reduced-form relationships are not presumed invariant after an equilibrium-changing intervention.

\textbf{Optimization and welfare.} Optimization is over the stated feasible set. If attainment is not established, $\sup$ or $\inf$ and $\varepsilon$-optimal choices replace an asserted maximizer or minimizer. For example, with value $V=\sup_{a\in\mathcal A}W(a)<\infty$, an $\varepsilon$-optimal set is $\{a\in\mathcal A:W(a)\geq V-\varepsilon\}$ for $\varepsilon>0$. Benefits, costs, risk and health or environmental outcomes can be aggregated only using declared units and conversion weights. Interpersonal utility weights, the treatment of fiscal transfers and foreign profits, discounting, and the behavioral welfare criterion are explicit inputs. A comparative-static derivative additionally requires existence and appropriate differentiability of the selected counterfactual.

\textbf{Sources.} Local labels [S1], [S2], and so forth restart in every question. Locators identify the relevant abstract, model, section or result at the level actually checked; a bibliographic or abstract-level verification is not represented as inspection of an entire proof. Working papers are distinguished from later publications and changing versions. Source summaries are paraphrased. All URL checks and editorial formulations refer to the audit date stated above.
% END ECONOMICS STATEMENT
\end{document}
